Understanding the role of dark matter in compact stars is a fascinating challenge in modern astrophysics. This study explores how dark matter influences the structure and properties of neutron stars, boson stars, and hybrid two-fluid stellar models. By solving the Tolman-Oppenheimer-Volkoff (TOV) equations for neutron stars and the Einstein-Klein-Gordon (EKG) equations for boson stars, we simulate their mass-radius relationships and tidal deformability. Using advanced computational tools in Python, MATLAB, and OCTAVE, we analyze how different potential interactions shape these stars' internal structures. Additionally, we investigate two-fluid models that incorporate interacting bosonic fields, allowing us to estimate the fraction of dark matter within compact stars. By comparing our results with gravitational wave (GW) observations, we provide new constraints on the presence of dark matter in stellar environments. This research not only deepens our understanding of compact objects but also offers valuable insights into the elusive nature of dark matter in the universe.